
\documentclass{article}
\usepackage{colortbl}
\usepackage{makecell}
\usepackage{multirow}
\usepackage{supertabular}

\begin{document}

\newcounter{utterance}

\centering \large Interaction Transcript for game `cryptolect', experiment `unguided', episode 9 with qwen3.5{-}9b.
\vspace{24pt}

{ \footnotesize  \setcounter{utterance}{1}
\setlength{\tabcolsep}{0pt}
\begin{supertabular}{c@{$\;$}|p{.15\linewidth}@{}p{.15\linewidth}p{.15\linewidth}p{.15\linewidth}p{.15\linewidth}p{.15\linewidth}}
    \# & \multicolumn{2}{c}{Player} && \multicolumn{2}{c}{Game Master} \\
    \hline

    \theutterance \stepcounter{utterance}  
    & & & \multicolumn{4}{p{0.6\linewidth}}{
        \cellcolor[rgb]{0.9,0.9,0.9}{
            \makecell[{{p{\linewidth}}}]{
                \texttt{\tiny{[P1$\langle$GM]}}
                \texttt{You are a field linguist. You have made contact with the only speaker of an unknown} \\
\texttt{language (the "informant"). Nobody has ever written down this language before, and you} \\
\texttt{have NO examples to start from. Your task is to work out how it works entirely on your} \\
\texttt{own, and then translate sentences into it correctly.} \\
\\ 
\texttt{You will study the language by asking the informant to translate English sentences for} \\
\texttt{you. You may only talk about a small fixed world, using exactly these English words:} \\
\\ 
\texttt{Nouns: tree, stone, bird, ship, road, hill, dog, spear, boat, river, hunter, king, farmer, lake} \\
\texttt{Adjectives: old, blue, dark, big, red, small, tall, cold} \\
\texttt{Verbs (write them in this base form): find, hold, chase, feed, watch, carry, see, hear} \\
\texttt{Articles: the (definite), a (indefinite)} \\
\texttt{Number words: two, three, four, five, six (no number word = one / singular)} \\
\texttt{These are ALL the words in play {-} the language has no other words.} \\
\\ 
\texttt{HOW TO WRITE AN ENGLISH SENTENCE} \\
\texttt{Every sentence is one clause with a subject, a verb, and an object, written in this fixed English notation:} \\
\texttt{    SUBJECT  verb  OBJECT} \\
\texttt{Each of SUBJECT and OBJECT is a noun phrase of the form:} \\
\texttt{    {[}the | a{]}  {[}number word{]}  {[}adjective ...{]}  noun} \\
\texttt{{-} 'the' = definite, 'a' = indefinite (singular only); leave the article off for an indefinite plural.} \\
\texttt{{-} a number word (two, three, four, five, six) sets the count; with no number word the noun is singular.} \\
\texttt{{-} zero or more adjectives may come before the noun.} \\
\texttt{Examples of well{-}formed English sentences you may send:} \\
\texttt{    a tree find a stone} \\
\texttt{    six trees find a stone} \\
\texttt{    a old tree find a stone} \\
\\ 
\texttt{HOW THE SESSION WORKS} \\
\texttt{1. INVESTIGATION: You have 5 rounds of questioning. In each round you may ask} \\
\texttt{   the informant to translate up to 4 English sentences (so at most} \\
\texttt{   20 probes in total). The informant answers each one truthfully and} \\
\texttt{   exactly. You are starting from nothing, so spend your early probes establishing the} \\
\texttt{   basic words and word order, then test how number, case, definiteness, agreement and} \\
\texttt{   any sound changes work. Choose your probes deliberately {-} test specific hypotheses.} \\
\texttt{2. EXAM: After the investigation rounds, the informant will give you 10 NEW} \\
\texttt{   English sentences. You must translate each one INTO the language. Every answer is} \\
\texttt{   graded by EXACT match against how the informant would say it {-} every word, ending and} \\
\texttt{   space must be right (capitalisation and surrounding spaces are ignored). The exam} \\
\texttt{   sentences will use the same vocabulary but in combinations you have not been shown, so} \\
\texttt{   you must induce the general rules, not just memorise examples.} \\
\\ 
\texttt{RESPONSE FORMAT (for every investigation round)} \\
\texttt{List your probe sentences under a 'PROBES:' line, one per line. A THOUGHTS line is} \\
\texttt{optional {-} keep any reasoning brief so you always finish writing your PROBES:} \\
\texttt{THOUGHTS: <optional; brief>} \\
\texttt{PROBES:} \\
\texttt{1. <an English sentence, written as described above>} \\
\texttt{2. <another English sentence>} \\
\texttt{(up to 4 lines; at least one)} \\
\\ 
\texttt{Begin your first round of probes now.} \\
            }
        }
    }
     \\ \\

    \theutterance \stepcounter{utterance}  
    & \multicolumn{4}{p{0.6\linewidth}}{
        \cellcolor[rgb]{0.561,0.737,0.561}{
            \makecell[{{p{\linewidth}}}]{
                \texttt{\tiny{[P1$\rangle$GM]}}
                \texttt{THOUGHTS: I need to establish the basic vocabulary and word order immediately. I will start with simple Subject{-}Verb{-}Object sentences using singular nouns and verbs to see if the order is SVO (most common) or something else. I'll use "tree", "find", and "stone" as my first set. I will also test the indefinite article "a" to see if it translates to a specific word or is omitted.} \\
\\ 
\texttt{PROBES:} \\
\texttt{1. a tree find a stone} \\
\texttt{2. a tree find stone} \\
\texttt{3. tree find a stone} \\
\texttt{4. tree find stone} \\
            }
        }
    }
    & & \\ \\

    \theutterance \stepcounter{utterance}  
    & & & \multicolumn{2}{p{0.3\linewidth}}{
        \cellcolor[rgb]{0.9,0.9,0.9}{
            \makecell[{{p{\linewidth}}}]{
                \texttt{\tiny{[GM$|$GM]}}
                {'kind': 'investigation', 'round': 0, 'probes_issued': 4, 'malformed': 0, 'probes': ['a tree find a stone', 'a tree find stone', 'tree find a stone', 'tree find stone']}
            }
        }
    }
    & & \\ \\

    \theutterance \stepcounter{utterance}  
    & & & \multicolumn{4}{p{0.6\linewidth}}{
        \cellcolor[rgb]{0.9,0.9,0.9}{
            \makecell[{{p{\linewidth}}}]{
                \texttt{\tiny{[P1$\langle$GM]}}
                \texttt{TRANSLATIONS:} \\
\texttt{1. "a tree find a stone"  =>  hobofi reho noba} \\
\texttt{2. "a tree find stone"  =>  hobofi reho noba} \\
\texttt{3. "tree find a stone"  =>  hobofi reho noba} \\
\texttt{4. "tree find stone"  =>  hobofi reho noba} \\
\texttt{STATUS: investigation round 1 of 5 done. Probes used: 4 of 20. 4 investigation round(s) left before the exam (10 sentences).} \\
\texttt{Send your next probes now (a PROBES list; brief THOUGHTS optional).} \\
            }
        }
    }
     \\ \\

    \theutterance \stepcounter{utterance}  
    & \multicolumn{4}{p{0.6\linewidth}}{
        \cellcolor[rgb]{0.561,0.737,0.561}{
            \makecell[{{p{\linewidth}}}]{
                \texttt{\tiny{[P1$\rangle$GM]}}
                \texttt{THOUGHTS: The results show that "a tree find a stone", "a tree find stone", "tree find a stone", and "tree find stone" all translate to the exact same string: "hobofi reho noba". This is a critical finding. It suggests that in this language, the indefinite article "a" and the definite article "the" are either non{-}existent, phonologically null, or the nouns/verbs do not change based on definiteness in this context. Furthermore, the lack of difference between singular and plural (if implied by context not present here) or the specific handling of "a" vs no "a" is confusing. Wait, looking closer at the prompt's grammar rules: "a number word = one / singular". So "a tree" is singular. "tree" alone is also singular. The fact that they are identical suggests the language might not mark number or definiteness on the noun phrase itself in a way that changes the word form, OR the translation provided by the informant is a fixed phrase for this specific semantic content regardless of the English syntactic variations I introduced.} \\
\\ 
\texttt{However, a more likely linguistic hypothesis given the data is that the language simply does not have articles, and the number "one" is not marked on the noun. The variations I tested (adding/removing 'a') produced no change. This implies the core structure is likely SVO (Subject{-}Verb{-}Object) based on the order: Noun1 (tree) {-} Verb (find) {-} Noun2 (stone).} \\
\\ 
\texttt{Let's test the verb "hold" and introduce a second noun to see if word order changes or if there are other particles. I also need to check if the language marks plurality explicitly with a suffix or a separate word, since "a tree" and "tree" were identical. I will try a sentence with two objects or a different verb to see if the pattern holds. Actually, let's test a different verb and a different noun pair to confirm the SVO order and see if the words change. I'll use "bird", "chase", "dog".} \\
\\ 
\texttt{PROBES:} \\
\texttt{1. a bird chase a dog} \\
\texttt{2. bird chase dog} \\
\texttt{3. a big tree hold a small stone} \\
\texttt{4. big tree hold small stone} \\
            }
        }
    }
    & & \\ \\

    \theutterance \stepcounter{utterance}  
    & & & \multicolumn{2}{p{0.3\linewidth}}{
        \cellcolor[rgb]{0.9,0.9,0.9}{
            \makecell[{{p{\linewidth}}}]{
                \texttt{\tiny{[GM$|$GM]}}
                {'kind': 'investigation', 'round': 1, 'probes_issued': 4, 'malformed': 0, 'probes': ['a bird chase a dog', 'bird chase dog', 'a big tree hold a small stone', 'big tree hold small stone']}
            }
        }
    }
    & & \\ \\

    \theutterance \stepcounter{utterance}  
    & & & \multicolumn{4}{p{0.6\linewidth}}{
        \cellcolor[rgb]{0.9,0.9,0.9}{
            \makecell[{{p{\linewidth}}}]{
                \texttt{\tiny{[P1$\langle$GM]}}
                \texttt{TRANSLATIONS:} \\
\texttt{1. "a bird chase a dog"  =>  hede rinera gubo} \\
\texttt{2. "bird chase dog"  =>  hede rinera gubo} \\
\texttt{3. "a big tree hold a small stone"  =>  hobofi nafegu fohu noba bina} \\
\texttt{4. "big tree hold small stone"  =>  hobofi nafegu fohu noba bina} \\
\texttt{STATUS: investigation round 2 of 5 done. Probes used: 8 of 20. 3 investigation round(s) left before the exam (10 sentences).} \\
\texttt{Send your next probes now (a PROBES list; brief THOUGHTS optional).} \\
            }
        }
    }
     \\ \\

    \theutterance \stepcounter{utterance}  
    & \multicolumn{4}{p{0.6\linewidth}}{
        \cellcolor[rgb]{0.561,0.737,0.561}{
            \makecell[{{p{\linewidth}}}]{
                \texttt{\tiny{[P1$\rangle$GM]}}
                \texttt{THOUGHTS: The data confirms that the English articles "a" and "the" are completely ignored or non{-}existent in this language's grammar for these sentences. The translations for "a tree find a stone" and "tree find stone" are identical, as are the variations with adjectives. This suggests the language does not mark definiteness or number (singular/plural) on the nouns themselves in the way English does, or perhaps the "number word" rule implies that without a number word >1, it defaults to singular, and "a" is just a semantic marker that doesn't translate.} \\
\\ 
\texttt{Crucially, I need to understand how adjectives work. In sentence 3 ("a big tree hold a small stone"), the translation is "hobofi nafegu fohu noba bina".} \\
\texttt{Comparing to sentence 1 ("a tree find a stone" {-}> "hobofi reho noba"):} \\
\texttt{{-} "tree" seems to be "hobofi".} \\
\texttt{{-} "find" seems to be "reho".} \\
\texttt{{-} "stone" seems to be "noba".} \\
\texttt{{-} "big" seems to be "nafegu".} \\
\texttt{{-} "hold" seems to be "fohu".} \\
\texttt{{-} "small" seems to be "bina".} \\
\\ 
\texttt{The order in sentence 3 is: {[}tree{]} {[}big{]} {[}hold{]} {[}stone{]} {[}small{]}.} \\
\texttt{Wait, let's re{-}parse "hobofi nafegu fohu noba bina".} \\
\texttt{If "hobofi" = tree, "nafegu" = big, "fohu" = hold, "noba" = stone, "bina" = small.} \\
\texttt{The order is: Noun1 Adj1 Verb Noun2 Adj2.} \\
\texttt{This is a very specific word order: S (Noun) {-} Adj {-} V {-} O (Noun) {-} Adj.} \\
\texttt{Let's verify this hypothesis with a new sentence where the adjective position matters or where I can test a different adjective/noun combination. I should also test if the adjective comes before or after the noun in a different context, or if the order is strictly N{-}Adj{-}V{-}O{-}Adj.} \\
\texttt{Actually, looking at "hobofi nafegu fohu noba bina":} \\
\texttt{Subject: a big tree {-}> hobofi nafegu} \\
\texttt{Verb: hold {-}> fohu} \\
\texttt{Object: a small stone {-}> noba bina} \\
\texttt{It looks like the adjective follows the noun? Or precedes?} \\
\texttt{In English: "big tree" (Adj{-}Noun). In translation: "hobofi nafegu" (Noun{-}Adj).} \\
\texttt{In English: "small stone" (Adj{-}Noun). In translation: "noba bina" (Noun{-}Adj).} \\
\texttt{So the rule seems to be: Noun comes first, then Adjective.} \\
\texttt{Let's test this with a new sentence to confirm the N{-}Adj order and check if the verb position is strictly between the two noun phrases. I will use "blue", "river", "boat", "see".} \\
\\ 
\texttt{PROBES:} \\
\texttt{1. a blue river see a boat} \\
\texttt{2. blue river see boat} \\
\texttt{3. a tall hunter carry a red spear} \\
\texttt{4. tall hunter carry red spear} \\
            }
        }
    }
    & & \\ \\

    \theutterance \stepcounter{utterance}  
    & & & \multicolumn{2}{p{0.3\linewidth}}{
        \cellcolor[rgb]{0.9,0.9,0.9}{
            \makecell[{{p{\linewidth}}}]{
                \texttt{\tiny{[GM$|$GM]}}
                {'kind': 'investigation', 'round': 2, 'probes_issued': 4, 'malformed': 0, 'probes': ['a blue river see a boat', 'blue river see boat', 'a tall hunter carry a red spear', 'tall hunter carry red spear']}
            }
        }
    }
    & & \\ \\

    \theutterance \stepcounter{utterance}  
    & & & \multicolumn{4}{p{0.6\linewidth}}{
        \cellcolor[rgb]{0.9,0.9,0.9}{
            \makecell[{{p{\linewidth}}}]{
                \texttt{\tiny{[P1$\langle$GM]}}
                \texttt{TRANSLATIONS:} \\
\texttt{1. "a blue river see a boat"  =>  nadu nifafu daba duha} \\
\texttt{2. "blue river see boat"  =>  nadu nifafu daba duha} \\
\texttt{3. "a tall hunter carry a red spear"  =>  feri hunege hebo nife gabega} \\
\texttt{4. "tall hunter carry red spear"  =>  feri hunege hebo nife gabega} \\
\texttt{STATUS: investigation round 3 of 5 done. Probes used: 12 of 20. 2 investigation round(s) left before the exam (10 sentences).} \\
\texttt{Send your next probes now (a PROBES list; brief THOUGHTS optional).} \\
            }
        }
    }
     \\ \\

    \theutterance \stepcounter{utterance}  
    & \multicolumn{4}{p{0.6\linewidth}}{
        \cellcolor[rgb]{0.561,0.737,0.561}{
            \makecell[{{p{\linewidth}}}]{
                \texttt{\tiny{[P1$\rangle$GM]}}
                \texttt{THOUGHTS: The data strongly supports the hypothesis that the language word order is Subject{-}Noun{-}Adjective{-}Verb{-}Object{-}Noun{-}Adjective (S{-}N{-}A{-}V{-}O{-}N{-}A).} \\
\texttt{Let's verify the components from the latest round:} \\
\texttt{1. "a blue river see a boat" {-}> "nadu nifafu daba duha"} \\
\texttt{   {-} Subject: blue river {-}> nadu nifafu (Noun: nadu, Adj: nifafu)} \\
\texttt{   {-} Verb: see {-}> daba} \\
\texttt{   {-} Object: boat {-}> duha (Wait, where is the adjective for boat? "a boat" has no adjective. So "duha" = boat).} \\
\texttt{   {-} Wait, let's re{-}examine sentence 1 vs 3.} \\
\texttt{   {-} Sentence 1: "blue river" (Adj{-}Noun in English) {-}> "nadu nifafu". So Noun=river=nadu, Adj=blue=nifafu.} \\
\texttt{   {-} Sentence 3: "tall hunter" {-}> "feri hunege". Noun=hunter=feri, Adj=tall=hunege.} \\
\texttt{   {-} Sentence 3: "red spear" {-}> "nife gabega". Noun=spear=gabega, Adj=red=nife.} \\
\texttt{   {-} Sentence 1: "boat" {-}> "duha".} \\
\texttt{   {-} Sentence 1: "see" {-}> "daba".} \\
\texttt{   {-} Structure: {[}Noun1{]} {[}Adj1{]} {[}Verb{]} {[}Noun2{]} {[}Adj2{]}?} \\
\texttt{   {-} Let's check Sentence 1 again: "nadu nifafu daba duha".} \\
\texttt{     {-} If Noun1=river=nadu, Adj1=blue=nifafu.} \\
\texttt{     {-} Verb=see=daba.} \\
\texttt{     {-} Noun2=boat=duha.} \\
\texttt{     {-} Where is the adjective for the object? There isn't one in "a boat".} \\
\texttt{     {-} But in Sentence 3: "feri hunege hebo nife gabega".} \\
\texttt{       {-} Noun1=hunter=feri, Adj1=tall=hunege.} \\
\texttt{       {-} Verb=carry=hebo.} \\
\texttt{       {-} Noun2=spear=gabega, Adj2=red=nife.} \\
\texttt{     {-} Order: N1 {-} A1 {-} V {-} N2 {-} A2.} \\
\texttt{     {-} Let's re{-}check Sentence 1: "nadu nifafu daba duha".} \\
\texttt{       {-} N1 (river) {-} A1 (blue) {-} V (see) {-} N2 (boat).} \\
\texttt{       {-} This fits perfectly.} \\
\texttt{   {-} Let's re{-}check Sentence 3: "feri hunege hebo nife gabega".} \\
\texttt{       {-} N1 (hunter) {-} A1 (tall) {-} V (carry) {-} A2 (red) {-} N2 (spear).} \\
\texttt{       {-} Wait, in Sentence 3, "red" comes BEFORE "spear". "nife gabega".} \\
\texttt{       {-} In Sentence 1, "blue" comes AFTER "river"? No, "nadu nifafu". "nifafu" is after "nadu".} \\
\texttt{       {-} So the order is Noun{-}Adjective for both subject and object?} \\
\texttt{       {-} Let's look at Sentence 1 again: "nadu nifafu". Noun=river, Adj=blue. Order: N{-}A.} \\
\texttt{       {-} Sentence 3: "nife gabega". Adj=red, Noun=spear. Order: A{-}N.} \\
\texttt{       {-} Contradiction!} \\
\texttt{       {-} Let's re{-}evaluate the mapping.} \\
\texttt{       {-} Maybe "nifafu" is not "blue"?} \\
\texttt{       {-} Maybe "nadu" is not "river"?} \\
\texttt{       {-} Let's look at the previous round (Round 2): "hobofi nafegu fohu noba bina".} \\
\texttt{         {-} "big tree hold small stone".} \\
\texttt{         {-} "hobofi nafegu": If N{-}A, then tree=hobofi, big=nafegu.} \\
\texttt{         {-} "noba bina": If A{-}N, then small=bina, stone=noba.} \\
\texttt{         {-} This would mean Subject is N{-}A and Object is A{-}N. That seems inconsistent.} \\
\texttt{       {-} Alternative: Maybe the adjective position depends on the noun? Unlikely for a simple language.} \\
\texttt{       {-} Alternative: Maybe I misidentified the words.} \\
\texttt{       {-} Let's look at Round 2 again: "hobofi nafegu fohu noba bina".} \\
\texttt{         {-} Could it be: {[}Noun{]} {[}Verb{]} {[}Noun{]} {[}Adj{]}? No, "hold" is clearly the verb.} \\
\texttt{         {-} Could it be: {[}Noun{]} {[}Adj{]} {[}Verb{]} {[}Noun{]} {[}Adj{]}?} \\
\texttt{           {-} hobofi (tree) nafegu (big) fohu (hold) noba (stone) bina (small).} \\
\texttt{           {-} Order: N{-}A{-}V{-}N{-}A.} \\
\texttt{         {-} Round 3 Sentence 1: "nadu nifafu daba duha".} \\
\texttt{           {-} nadu (river) nifafu (blue) daba (see) duha (boat).} \\
\texttt{           {-} Order: N{-}A{-}V{-}N. (No object adjective).} \\
\texttt{           {-} This matches N{-}A{-}V{-}N{-}A if the last part is missing.} \\
\texttt{         {-} Round 3 Sentence 3: "feri hunege hebo nife gabega".} \\
\texttt{           {-} feri (hunter) hunege (tall) hebo (carry) nife (red) gabega (spear).} \\
\texttt{           {-} Order: N{-}A{-}V{-}A{-}N.} \\
\texttt{           {-} Here the object adjective comes BEFORE the noun.} \\
\texttt{       {-} This is a major inconsistency. N{-}A for subject, A{-}N for object?} \\
\texttt{       {-} Let's reconsider the word boundaries.} \\
\texttt{       {-} Is it possible "nife" is the noun and "gabega" is the adjective? "red spear".} \\
\texttt{       {-} Is it possible "noba" is the adjective and "bina" is the noun? "small stone".} \\
\texttt{       {-} If "noba" = small and "bina" = stone.} \\
\texttt{       {-} Then Round 2 Object: "noba bina" = small stone (A{-}N).} \\
\texttt{       {-} Round 3 Object: "nife gabega" = red spear (A{-}N).} \\
\texttt{       {-} Round 2 Subject: "hobofi nafegu". If N{-}A, then tree=hobofi, big=nafegu.} \\
\texttt{       {-} Round 3 Subject: "nadu nifafu". If N{-}A, then river=nadu, blue=nifafu.} \\
\texttt{       {-} Round 3 Subject 2: "feri hunege". If N{-}A, then hunter=feri, tall=hunege.} \\
\texttt{       {-} So Subject is always N{-}A. Object is always A{-}N?} \\
\texttt{       {-} Or maybe the language has a specific rule: Adjectives follow the subject noun but precede the object noun?} \\
\texttt{       {-} Let's test this hypothesis: **Subject{-}Noun{-}Adjective{-}Verb{-}Object{-}Noun{-}Adjective**? No, that would be N{-}A{-}V{-}N{-}A.} \\
\texttt{       {-} My analysis of Round 3 Obj: "nife gabega" (red spear). If A{-}N, then "nife"=red, "gabega"=spear.} \\
\texttt{       {-} My analysis of Round 2 Obj: "noba bina" (small stone). If A{-}N, then "noba"=small, "bina"=stone.} \\
\texttt{       {-} My analysis of Round 2 Subj: "hobofi nafegu" (big tree). If N{-}A, then "hobofi"=tree, "nafegu"=big.} \\
\texttt{       {-} My analysis of Round 3 Subj: "nadu nifafu" (blue river). If N{-}A, then "nadu"=river, "nifafu"=blue.} \\
\texttt{       {-} My analysis of Round 3 Subj: "feri hunege" (tall hunter). If N{-}A, then "feri"=hunter, "hunege"=tall.} \\
\texttt{       {-} Hypothesis: **Adjectives follow Subject Nouns but precede Object Nouns.**} \\
\texttt{       {-} Let's try to verify this with a new sentence where I can swap the adjective position or use a different structure.} \\
\texttt{       {-} I need to confirm if the order is strictly N{-}A for Subject and A{-}N for Object.} \\
\texttt{       {-} I will construct a sentence with a Subject that has an adjective and an Object that has an adjective, and see if the pattern holds.} \\
\texttt{       {-} I will also test a sentence with NO adjectives to see the base SVO order.} \\
\texttt{       {-} And I will test a sentence with TWO adjectives on the subject to see if they stack (N{-}A1{-}A2 or N{-}A2{-}A1).} \\
\texttt{       {-} Let's try: "a big blue tree find a small red stone".} \\
\texttt{       {-} If Subject is N{-}A1{-}A2: tree big blue {-}> hobofi nafegu nifafu?} \\
\texttt{       {-} If Object is A1{-}N{-}A2: stone small red {-}> noba bina nife?} \\
\texttt{       {-} This will clarify the stacking order and the subject/object distinction.} \\
\\ 
\texttt{PROBES:} \\
\texttt{1. a big blue tree find a small red stone} \\
\texttt{2. big blue tree find small red stone} \\
\texttt{3. a tree find a stone} \\
\texttt{4. a big tree find a big stone} \\
            }
        }
    }
    & & \\ \\

    \theutterance \stepcounter{utterance}  
    & & & \multicolumn{2}{p{0.3\linewidth}}{
        \cellcolor[rgb]{0.9,0.9,0.9}{
            \makecell[{{p{\linewidth}}}]{
                \texttt{\tiny{[GM$|$GM]}}
                {'kind': 'investigation', 'round': 3, 'probes_issued': 4, 'malformed': 0, 'probes': ['a big blue tree find a small red stone', 'big blue tree find small red stone', 'a tree find a stone', 'a big tree find a big stone']}
            }
        }
    }
    & & \\ \\

    \theutterance \stepcounter{utterance}  
    & & & \multicolumn{4}{p{0.6\linewidth}}{
        \cellcolor[rgb]{0.9,0.9,0.9}{
            \makecell[{{p{\linewidth}}}]{
                \texttt{\tiny{[P1$\langle$GM]}}
                \texttt{TRANSLATIONS:} \\
\texttt{1. "a big blue tree find a small red stone"  =>  hobofi nafegu nifafu reho noba bina gabega} \\
\texttt{2. "big blue tree find small red stone"  =>  hobofi nafegu nifafu reho noba bina gabega} \\
\texttt{3. "a tree find a stone"  =>  hobofi reho noba} \\
\texttt{4. "a big tree find a big stone"  =>  hobofi nafegu reho noba nafegu} \\
\texttt{STATUS: investigation round 4 of 5 done. Probes used: 16 of 20. 1 investigation round(s) left before the exam (10 sentences).} \\
\texttt{Send your next probes now (a PROBES list; brief THOUGHTS optional).} \\
            }
        }
    }
     \\ \\

    \theutterance \stepcounter{utterance}  
    & \multicolumn{4}{p{0.6\linewidth}}{
        \cellcolor[rgb]{0.561,0.737,0.561}{
            \makecell[{{p{\linewidth}}}]{
                \texttt{\tiny{[P1$\rangle$GM]}}
                \texttt{THOUGHTS: The data confirms my hypothesis about adjective placement: Adjectives follow Subject Nouns (N{-}A) but precede Object Nouns (A{-}N).} \\
\texttt{{-} Round 4, Sentence 1: "big blue tree" {-}> "hobofi nafegu nifafu" (N{-}A1{-}A2). "small red stone" {-}> "noba bina gabega" (N{-}A1{-}A2? Wait).} \\
\texttt{Let's re{-}examine Round 4, Sentence 1: "hobofi nafegu nifafu reho noba bina gabega".} \\
\texttt{Subject: "big blue tree". Translation: "hobofi nafegu nifafu".} \\
\texttt{Order: Noun (hobofi) {-} Adj1 (nafegu/big) {-} Adj2 (nifafu/blue). So Subject is N{-}A{-}A.} \\
\texttt{Object: "small red stone". Translation: "noba bina gabega".} \\
\texttt{Order: Noun (noba) {-} Adj1 (bina/small) {-} Adj2 (gabega/red).} \\
\texttt{Wait, in Round 3, "red spear" was "nife gabega" (A{-}N). Here "red stone" is "gabega" (part of N{-}A{-}A?).} \\
\texttt{Let's re{-}evaluate the word "gabega".} \\
\texttt{Round 3: "red spear" {-}> "nife gabega". If A{-}N, then nife=red, gabega=spear.} \\
\texttt{Round 4: "red stone" {-}> "gabega" is the last word. If the pattern is N{-}A{-}A for objects, then "noba" (stone) {-} "bina" (small) {-} "gabega" (red).} \\
\texttt{This implies "gabega" = red.} \\
\texttt{But in Round 3, "nife" was the first word of the object phrase "nife gabega". If "gabega" is red, then "nife" must be "spear".} \\
\texttt{So Round 3 "red spear" = "nife (spear) gabega (red)". This is A{-}N order? No, that would be Noun{-}Adj.} \\
\texttt{Let's re{-}read Round 3 carefully: "a tall hunter carry a red spear" {-}> "feri hunege hebo nife gabega".} \\
\texttt{Subject: "tall hunter" {-}> "feri hunege". (Noun=feri, Adj=hunege). Order N{-}A.} \\
\texttt{Object: "red spear". Translation: "nife gabega".} \\
\texttt{If "gabega" = red (from Round 4), then "nife" = spear.} \\
\texttt{So Object is "spear red" (N{-}A).} \\
\texttt{But in Round 2, "small stone" was "noba bina". If "noba"=stone, "bina"=small. Order N{-}A.} \\
\texttt{In Round 4, "small red stone" is "noba bina gabega". Order N{-}A{-}A.} \\
\texttt{So the rule is consistent: **Subject and Object both follow Noun{-}Adjective order (N{-}A...)**.} \\
\texttt{My previous confusion about "A{-}N" for objects was incorrect. "nife gabega" in Round 3 is likely "spear red" (N{-}A), meaning "nife"=spear, "gabega"=red.} \\
\texttt{Let's verify "nife" vs "noba".} \\
\texttt{Round 2: "stone" = noba.} \\
\texttt{Round 3: "spear" = nife.} \\
\texttt{Round 4: "red" = gabega.} \\
\texttt{So "red spear" = "nife gabega" (Spear Red). Correct.} \\
\texttt{"small stone" = "noba bina" (Stone Small). Correct.} \\
\texttt{"small red stone" = "noba bina gabega" (Stone Small Red). Correct.} \\
\texttt{"big tree" = "hobofi nafegu" (Tree Big). Correct.} \\
\texttt{"big blue tree" = "hobofi nafegu nifafu" (Tree Big Blue). Correct.} \\
\\ 
\texttt{So the grammar is:} \\
\texttt{1. Word Order: Subject{-}Verb{-}Object (SVO).} \\
\texttt{2. Noun Phrase Structure: Noun followed by Adjectives (N{-}A1{-}A2...).} \\
\texttt{3. Articles/Numbers: "a" and "the" are ignored. Number words >1 are not tested yet, but the prompt says "no number word = one". I should test plural explicitly if possible, but I have 1 probe left.} \\
\texttt{4. I need to check if number words (two, three, etc.) change the noun or adjective order, or if they just sit in a specific slot.} \\
\texttt{5. I also need to check if the verb changes form (conjugation) based on subject/object.} \\
\texttt{6. I have 1 probe left. I should test a number word to see where it goes.} \\
\texttt{   Hypothesis: {[}Noun{]} {[}Number{]} {[}Adj{]} ... or {[}Number{]} {[}Noun{]} ...} \\
\texttt{   Let's try: "two big trees find a stone".} \\
\texttt{   If number is separate: "hobofi nafegu hobofi"? No, "two trees".} \\
\texttt{   Maybe "hobofi hobofi"? No.} \\
\texttt{   Let's try "two trees find a stone".} \\
\texttt{   If "two" is a separate word, where does it go?} \\
\texttt{   Common patterns: Num{-}Noun or Noun{-}Num.} \\
\texttt{   Let's try: "two trees find a stone".} \\
\texttt{   Also, I should check if the verb changes. "find" (reho) vs "hold" (fohu). No obvious change yet.} \\
\texttt{   Let's test "two trees find a stone".} \\
\texttt{   And maybe "three birds chase a dog" to see if the pattern holds for other nouns.} \\
\texttt{   Actually, I only have 1 probe. I should maximize information.} \\
\texttt{   "two big trees find a small stone".} \\
\texttt{   This tests:} \\
\texttt{   {-} Number word placement.} \\
\texttt{   {-} Adjective placement with numbers.} \\
\texttt{   {-} Plural marking (if any).} \\
\texttt{   {-} Verb consistency.} \\
\\ 
\texttt{PROBES:} \\
\texttt{1. two big trees find a small stone} \\
            }
        }
    }
    & & \\ \\

    \theutterance \stepcounter{utterance}  
    & & & \multicolumn{2}{p{0.3\linewidth}}{
        \cellcolor[rgb]{0.9,0.9,0.9}{
            \makecell[{{p{\linewidth}}}]{
                \texttt{\tiny{[GM$|$GM]}}
                {'kind': 'investigation', 'round': 4, 'probes_issued': 1, 'malformed': 0, 'probes': ['two big trees find a small stone']}
            }
        }
    }
    & & \\ \\

    \theutterance \stepcounter{utterance}  
    & & & \multicolumn{4}{p{0.6\linewidth}}{
        \cellcolor[rgb]{0.9,0.9,0.9}{
            \makecell[{{p{\linewidth}}}]{
                \texttt{\tiny{[P1$\langle$GM]}}
                \texttt{TRANSLATIONS:} \\
\texttt{1. "two big trees find a small stone"  =>  hobofiru nafegu ni reho noba bina} \\
\\ 
\texttt{FINAL EXAM. Investigation is over. Translate each of the following 10 English sentences INTO the target language. Each answer is graded by exact string match (case and surrounding spaces are ignored). Every sentence uses only words you have already seen translated.} \\
\\ 
\texttt{1. the four tall roads find the four cold farmers} \\
\texttt{2. the two birds watch a old lake} \\
\texttt{3. two tall birds chase the four dark hunters} \\
\texttt{4. the four small lakes hear two red hunters} \\
\texttt{5. two tall lakes watch four dark spears} \\
\texttt{6. the cold tree chase two rivers} \\
\texttt{7. two tall ships carry two blue birds} \\
\texttt{8. five blue rivers see three small lakes} \\
\texttt{9. the spear hold six hills} \\
\texttt{10. the two lakes chase the cold dog} \\
\\ 
\texttt{Respond with the answers (you may add a brief THOUGHTS line first):} \\
\texttt{ANSWERS:} \\
\texttt{1. <translation of sentence 1>} \\
\texttt{... through 10.} \\
            }
        }
    }
     \\ \\

    \theutterance \stepcounter{utterance}  
    & \multicolumn{4}{p{0.6\linewidth}}{
        \cellcolor[rgb]{0.561,0.737,0.561}{
            \makecell[{{p{\linewidth}}}]{
                \texttt{\tiny{[P1$\rangle$GM]}}
                \texttt{THOUGHTS: Based on the investigation, I have deduced the following grammar rules:} \\
\texttt{1.  **Word Order**: Subject {-} Verb {-} Object (SVO).} \\
\texttt{2.  **Noun Phrase Structure**: Noun is followed by Adjectives (N{-}A1{-}A2...).} \\
\texttt{3.  **Number Placement**: The number word comes immediately after the Noun (N{-}Num{-}A...).} \\
\texttt{    *   Evidence: "two big trees" {-}> "hobofiru nafegu ni".} \\
\texttt{        *   "hobofi" = tree.} \\
\texttt{        *   "ru" = suffix for "two"? Or "hobofiru" is the word for "two trees"?} \\
\texttt{        *   Let's re{-}examine "hobofi reho noba" (tree find stone).} \\
\texttt{        *   "hobofiru nafegu ni" (two big trees...).} \\
\texttt{        *   If "hobofi" is tree, "hobofiru" looks like "tree" + "two" suffix? Or "two" + "tree"?} \\
\texttt{        *   Let's look at "four tall roads". "four" is "ni"? No, "ni" was in "hobofiru nafegu ni".} \\
\texttt{        *   Wait, "hobofiru nafegu ni".} \\
\texttt{        *   "hobofi" = tree.} \\
\texttt{        *   "nafegu" = big.} \\
\texttt{        *   "ni" = ?} \\
\texttt{        *   In "two big trees", "ni" is at the end of the subject phrase.} \\
\texttt{        *   In "four tall roads", "four" is the number.} \\
\texttt{        *   Let's look at the number words provided: two, three, four, five, six.} \\
\texttt{        *   In "hobofiru nafegu ni", if "hobofi" is tree, maybe "ru" is the plural marker for 2? And "ni" is... wait.} \\
\texttt{        *   Let's re{-}parse "hobofiru nafegu ni".} \\
\texttt{        *   Maybe "hobofiru" = two trees? And "nafegu" = big? And "ni" is extra? No.} \\
\texttt{        *   Maybe "hobofi" = tree. "ru" = two? "hobofiru" = two trees?} \\
\texttt{        *   Then "nafegu" = big.} \\
\texttt{        *   Then "ni" = ?} \\
\texttt{        *   Let's look at "four tall roads". "four" is a number.} \\
\texttt{        *   Let's look at "two birds".} \\
\texttt{        *   Let's reconsider the "two big trees" translation: "hobofiru nafegu ni".} \\
\texttt{        *   Could "ni" be the number "two"?} \\
\texttt{        *   If "ni" = two. Then "hobofiru" = big trees? No, "big" is "nafegu".} \\
\texttt{        *   Maybe "hobofi" = tree. "ru" = plural? "hobofiru" = trees (plural).} \\
\texttt{        *   Then "nafegu" = big.} \\
\texttt{        *   Then "ni" = two.} \\
\texttt{        *   Order: Noun{-}Plural{-}Adj{-}Num? That seems complex.} \\
\texttt{        *   Alternative: "hobofi" = tree. "ru" = two. "hobofiru" = two trees.} \\
\texttt{        *   Then "nafegu" = big.} \\
\texttt{        *   Then "ni" = ?} \\
\texttt{        *   Maybe "ni" is a particle?} \\
\texttt{        *   Let's look at "four tall roads". "four" is the number.} \\
\texttt{        *   Let's assume the number word is a separate word.} \\
\texttt{        *   "two big trees" {-}> "hobofiru nafegu ni".} \\
\texttt{        *   If "ni" = two. Then "hobofiru nafegu" = big trees.} \\
\texttt{        *   "hobofi" = tree. "ru" = plural suffix? "hobofiru" = trees.} \\
\texttt{        *   "nafegu" = big.} \\
\texttt{        *   So "trees big two"? No, SVO is fixed.} \\
\texttt{        *   Maybe the order is N{-}Num{-}A?} \\
\texttt{        *   "hobofi" (tree) "ni" (two) "nafegu" (big)? {-}> "hobofiru nafegu ni". No, "ni" is at the end.} \\
\texttt{        *   Maybe "ni" is not "two".} \\
\texttt{        *   Let's look at the list of numbers: two, three, four, five, six.} \\
\texttt{        *   Is it possible "ni" is "three"? No, "two" is the number.} \\
\texttt{        *   Let's try to map the numbers from the exam sentences to the probe sentences.} \\
\texttt{        *   Probe: "two big trees" {-}> "hobofiru nafegu ni".} \\
\texttt{        *   Exam 1: "four tall roads".} \\
\texttt{        *   Exam 2: "two birds".} \\
\texttt{        *   Exam 3: "four dark hunters".} \\
\texttt{        *   Exam 4: "two red hunters".} \\
\texttt{        *   Exam 5: "four dark spears".} \\
\texttt{        *   Exam 6: "two rivers".} \\
\texttt{        *   Exam 7: "two blue birds", "two blue birds".} \\
\texttt{        *   Exam 8: "five blue rivers", "three small lakes".} \\
\texttt{        *   Exam 9: "six hills".} \\
\texttt{        *   Exam 10: "two lakes", "cold dog".} \\
\\ 
\texttt{        Let's re{-}evaluate "hobofiru nafegu ni".} \\
\texttt{        Maybe "hobofi" = tree.} \\
\texttt{        "ru" = suffix for plural?} \\
\texttt{        "nafegu" = big.} \\
\texttt{        "ni" = two.} \\
\texttt{        Order: N{-}Plural{-}Adj{-}Num?} \\
\texttt{        Or maybe "hobofiru" = two trees. "nafegu" = big. "ni" = ?} \\
\texttt{        What if "ni" is the number "two"?} \\
\texttt{        Then "hobofiru nafegu" = big trees.} \\
\texttt{        "hobofi" = tree. "ru" = plural.} \\
\texttt{        So "trees big two".} \\
\texttt{        This implies the number comes LAST in the noun phrase?} \\
\texttt{        Let's check "four tall roads".} \\
\texttt{        If number is last: "roads tall four".} \\
\texttt{        Let's check "two birds". "birds two".} \\
\texttt{        Let's check "five blue rivers". "rivers blue five".} \\
\texttt{        Let's check "three small lakes". "lakes small three".} \\
\texttt{        Let's check "six hills". "hills six".} \\
\texttt{        This seems consistent. The number word is the last element of the noun phrase.} \\
\texttt{        So structure: {[}Noun{]} {[}Plural Marker?{]} {[}Adjectives{]} {[}Number{]}.} \\
\texttt{        Wait, "hobofiru". If "hobofi" is tree, "ru" is plural.} \\
\texttt{        Is "ru" always plural?} \\
\texttt{        "two big trees" {-}> "hobofiru".} \\
\texttt{        "four tall roads" {-}> "road" is "nife"? No, "spear" is "nife". "road" is unknown.} \\
\texttt{        Let's look at "four tall roads".} \\
\texttt{        If "ru" is plural, maybe "four" adds a different suffix?} \\
\texttt{        Or maybe "hobofiru" is just "two trees" and "ru" is part of the number word?} \\
\texttt{        No, "two" is a separate word in English.} \\
\texttt{        Let's assume the number word is a separate word at the end.} \\
\texttt{        "hobofi" = tree.} \\
\texttt{        "ru" = plural marker?} \\
\texttt{        "nafegu" = big.} \\
\texttt{        "ni" = two.} \\
\texttt{        So: Tree{-}Plural{-}Big{-}Two.} \\
\texttt{        Let's test this hypothesis on "four tall roads".} \\
\texttt{        "road" = ?} \\
\texttt{        "tall" = "hunege" (from "tall hunter").} \\
\texttt{        "four" = ?} \\
\texttt{        If "ni" = two, what is four?} \\
\texttt{        Maybe "ni" is not two.} \\
\texttt{        Let's look at "hobofiru nafegu ni" again.} \\
\texttt{        Could "ni" be "big"? No, "nafegu" is big.} \\
\texttt{        Could "ni" be "stone"? No, "noba" is stone.} \\
\texttt{        Could "ni" be "two"? Yes, likely.} \\
\texttt{        So "two" is "ni".} \\
\texttt{        Then "hobofiru" = big trees.} \\
\texttt{        "hobofi" = tree. "ru" = plural.} \\
\texttt{        So "tree{-}plural{-}big{-}two".} \\
\texttt{        Let's check "four tall roads".} \\
\texttt{        "road" = ?} \\
\texttt{        "tall" = "hunege".} \\
\texttt{        "four" = ?} \\
\texttt{        If "ni" = two, maybe "four" is "ni" + something? Or a different word.} \\
\texttt{        Let's look at "five blue rivers".} \\
\texttt{        "river" = "nadu" (from "blue river").} \\
\texttt{        "blue" = "nifafu".} \\
\texttt{        "five" = ?} \\
\texttt{        "rivers" = "nadu" + plural?} \\
\texttt{        If the pattern is N{-}Plural{-}Adj{-}Num.} \\
\texttt{        "nadu" + "ru" + "nifafu" + "five".} \\
\texttt{        Let's check "three small lakes".} \\
\texttt{        "lake" = ?} \\
\texttt{        "small" = "bina".} \\
\texttt{        "three" = ?} \\
\texttt{        "lakes" = "lake" + plural.} \\
\texttt{        Let's check "six hills".} \\
\texttt{        "hill" = ?} \\
\texttt{        "six" = ?} \\
\texttt{        "hills" = "hill" + plural.} \\
\\ 
\texttt{        Let's try to identify the number words from the exam sentences by looking for patterns or common roots.} \\
\texttt{        Actually, I can just construct the sentences based on the pattern:} \\
\texttt{        {[}Noun{]} {[}Plural Suffix?{]} {[}Adjectives{]} {[}Number Word{]}.} \\
\texttt{        Wait, in "hobofiru nafegu ni", "ru" is attached to "hobofi".} \\
\texttt{        Is "ru" the plural marker for ALL nouns?} \\
\texttt{        "hobofi" {-}> "hobofiru".} \\
\texttt{        "nadu" {-}> "naduru"?} \\
\texttt{        "nife" {-}> "niferu"?} \\
\texttt{        "feri" {-}> "feriru"?} \\
\texttt{        "hobofi" ends in 'i'. "hobofiru" ends in 'ru'.} \\
\texttt{        Maybe the plural marker is "{-}ru".} \\
\texttt{        Let's assume:} \\
\texttt{        Singular: Noun} \\
\texttt{        Plural: Noun + ru} \\
\texttt{        Adjectives follow the noun (and plural marker).} \\
\texttt{        Number word is at the very end.} \\
\texttt{        So: Noun{-}ru{-}Adj{-}Num.} \\
\texttt{        Let's verify "hobofiru nafegu ni".} \\
\texttt{        Tree{-}ru{-}Big{-}Two. Matches.} \\
\texttt{        Now I need the words for:} \\
\texttt{        {-} road} \\
\texttt{        {-} bird} \\
\texttt{        {-} hunter} \\
\texttt{        {-} lake} \\
\texttt{        {-} river} \\
\texttt{        {-} ship} \\
\texttt{        {-} hill} \\
\texttt{        {-} dog} \\
\texttt{        {-} spear} \\
\texttt{        {-} boat} \\
\texttt{        {-} king} \\
\texttt{        {-} farmer} \\
\texttt{        {-} find, watch, chase, hear, carry, see, hold (verbs)} \\
\texttt{        {-} adjectives: old, blue, dark, big, red, small, tall, cold} \\
\texttt{        {-} numbers: two (ni), three (?), four (?), five (?), six (?)} \\
\\ 
\texttt{        Let's deduce the number words.} \\
\texttt{        "two" = ni.} \\
\texttt{        "three" = ?} \\
\texttt{        "four" = ?} \\
\texttt{        "five" = ?} \\
\texttt{        "six" = ?} \\
\texttt{        I don't have direct translations for 3, 4, 5, 6.} \\
\texttt{        Wait, I must have missed something.} \\
\texttt{        Let's re{-}read the probe translations.} \\
\texttt{        1. "a tree find a stone" => hobofi reho noba} \\
\texttt{        2. "a bird chase a dog" => hede rinera gubo} \\
\texttt{        3. "a big tree hold a small stone" => hobofi nafegu fohu noba bina} \\
\texttt{        4. "a blue river see a boat" => nadu nifafu daba duha} \\
\texttt{        5. "a tall hunter carry a red spear" => feri hunege hebo nife gabega} \\
\texttt{        6. "two big trees find a small stone" => hobofiru nafegu ni reho noba bina} \\
\\ 
\texttt{        From 6:} \\
\texttt{        "hobofiru" = two big trees? Or big trees (plural) and "ni" = two?} \\
\texttt{        If "ni" = two.} \\
\texttt{        Then "hobofiru" = big trees.} \\
\texttt{        "hobofi" = tree. "ru" = plural.} \\
\texttt{        "nafegu" = big.} \\
\texttt{        "noba" = stone. "bina" = small.} \\
\texttt{        "reho" = find.} \\
\texttt{        So "two big trees find a small stone" = {[}Tree{-}ru{-}Big{]} {[}Two{]} {[}Find{]} {[}Stone{-}Plural? No, "a small stone" is singular{]}.} \\
\texttt{        Wait, "a small stone" is singular. "noba bina".} \\
\texttt{        So "ni" is definitely "two".} \\
\texttt{        And "ru" is the plural marker.} \\
\texttt{        So "hobofiru" = trees.} \\
\texttt{        "nafegu" = big.} \\
\texttt{        "ni" = two.} \\
\texttt{        Order: Noun{-}Plural{-}Adj{-}Num.} \\
\texttt{        Wait, "hobofiru nafegu ni".} \\
\texttt{        Tree{-}Plural{-}Big{-}Two.} \\
\texttt{        This works.} \\
\\ 
\texttt{        Now I need the words for the other nouns and numbers.} \\
\texttt{        I don't have the words for 3, 4, 5, 6.} \\
\texttt{        Is it possible the number words are derived? No, the prompt says "Number words: two, three, four, five, six".} \\
\texttt{        Maybe I can infer them from the exam sentences if there are clues?} \\
\texttt{        No, the exam sentences are new.} \\
\texttt{        Wait, did I miss a probe?} \\
\texttt{        I only did 4 probes per round.} \\
\texttt{        Maybe I can guess the number words?} \\
\texttt{        Or maybe the number words are not needed for the exam?} \\
\texttt{        Exam 1: "four tall roads". Need "four" and "road".} \\
\texttt{        Exam 2: "two birds". Need "bird" and "two" (ni).} \\
\texttt{        Exam 3: "four dark hunters". Need "four", "dark", "hunter".} \\
\texttt{        Exam 4: "two red hunters". Need "red", "hunter".} \\
\texttt{        Exam 5: "four dark spears". Need "four", "dark", "spear".} \\
\texttt{        Exam 6: "two rivers". Need "river", "two".} \\
\texttt{        Exam 7: "two tall ships". Need "ship", "tall".} \\
\texttt{        Exam 8: "five blue rivers", "three small lakes". Need "five", "three", "lake".} \\
\texttt{        Exam 9: "six hills". Need "six", "hill".} \\
\texttt{        Exam 10: "two lakes", "cold dog". Need "lake", "dog", "cold".} \\
\\ 
\texttt{        I am missing the words for:} \\
\texttt{        {-} road} \\
\texttt{        {-} bird} \\
\texttt{        {-} hunter} \\
\texttt{        {-} lake} \\
\texttt{        {-} river (nadu)} \\
\texttt{        {-} ship} \\
\texttt{        {-} hill} \\
\texttt{        {-} dog} \\
\texttt{        {-} spear (nife)} \\
\texttt{        {-} boat (duha)} \\
\texttt{        {-} king} \\
\texttt{        {-} farmer} \\
\texttt{        {-} verbs: watch, chase, hear, carry, see, hold, find (reho)} \\
\texttt{        {-} adjectives: old, dark, cold, red, small, tall, big, blue} \\
\texttt{        {-} numbers: three, four, five, six (two=ni)} \\
\\ 
\texttt{        This is a problem. I cannot translate without the words.} \\
\texttt{        Did I miss a clue?} \\
\texttt{        "hobofi" = tree.} \\
\texttt{        "noba" = stone.} \\
\texttt{        "hede" = bird? (from "a bird chase a dog" => hede rinera gubo).} \\
\texttt{        "rinera" = chase.} \\
\texttt{        "gubo" = dog.} \\
\texttt{        "nadu" = river.} \\
\texttt{        "nifafu" = blue.} \\
\texttt{        "daba" = see.} \\
\texttt{        "duha" = boat.} \\
\texttt{        "feri" = hunter.} \\
\texttt{        "hunege" = tall.} \\
\texttt{        "hebo" = carry.} \\
\texttt{        "nife" = spear.} \\
\texttt{        "gabega" = red.} \\
\texttt{        "hobofi" = tree.} \\
\texttt{        "nafegu" = big.} \\
\texttt{        "fohu" = hold.} \\
\texttt{        "bina" = small.} \\
\texttt{        "reho" = find.} \\
\\ 
\texttt{        Missing words:} \\
\texttt{        {-} road} \\
\texttt{        {-} lake} \\
\texttt{        {-} ship} \\
\texttt{        {-} hill} \\
\texttt{        {-} dog (gubo) {-} Found!} \\
\texttt{        {-} bird (hede) {-} Found!} \\
\texttt{        {-} hunter (feri) {-} Found!} \\
\texttt{        {-} river (nadu) {-} Found!} \\
\texttt{        {-} spear (nife) {-} Found!} \\
\texttt{        {-} boat (duha) {-} Found!} \\
\texttt{        {-} king} \\
\texttt{        {-} farmer} \\
\texttt{        {-} old} \\
\texttt{        {-} dark} \\
\texttt{        {-} cold} \\
\texttt{        {-} three} \\
\texttt{        {-} four} \\
\texttt{        {-} five} \\
\texttt{        {-} six} \\
\\ 
\texttt{        Wait, I have "gubo" = dog.} \\
\texttt{        I have "hede" = bird.} \\
\texttt{        I have "feri" = hunter.} \\
\texttt{        I have "nadu" = river.} \\
\texttt{        I have "nife" = spear.} \\
\texttt{        I have "duha" = boat.} \\
\texttt{        I have "hobofi" = tree.} \\
\texttt{        I have "noba" = stone.} \\
\texttt{        I have "hobofi" = tree.} \\
\texttt{        I have "noba" = stone.} \\
\\ 
\texttt{        What about "road", "lake", "ship", "hill", "king", "farmer"?} \\
\texttt{        And "old", "dark", "cold"?} \\
\texttt{        And "three", "four", "five", "six"?} \\
\texttt{        I must have missed something in the prompt or the translations.} \\
\texttt{        Let's re{-}read the translations carefully.} \\
\texttt{        Maybe "hobofiru" contains the number?} \\
\texttt{        "hobofiru" = two trees.} \\
\texttt{        "nafegu" = big.} \\
\texttt{        "ni" = ?} \\
\texttt{        If "ni" is not "two", what is it?} \\
\texttt{        Maybe "ni" is a particle?} \\
\texttt{        But "two" is a number word.} \\
\texttt{        If "ni" is not "two", then where is "two"?} \\
\texttt{        Maybe "hobofiru" = two trees.} \\
\texttt{        Then "nafegu" = big.} \\
\texttt{        Then "ni" = ?} \\
\texttt{        This doesn't make sense.} \\
\texttt{        Let's assume "ni" = two.} \\
\texttt{        Then I need the other numbers.} \\
\texttt{        Is it possible the number words are formed by adding a suffix to "ni"?} \\
\texttt{        ni (2), ni{-}? (3), ni{-}? (4)...} \\
\texttt{        Or maybe the number words are hidden in the translations I haven't fully parsed?} \\
\texttt{        No, the translations are fixed.} \\
\texttt{        Wait, I might have misidentified "hobofiru".} \\
\texttt{        "hobofi" = tree.} \\
\texttt{        "ru" = plural.} \\
\texttt{        "hobofiru" = trees.} \\
\texttt{        "nafegu" = big.} \\
\texttt{        "ni" = two.} \\
\texttt{        This seems the most logical.} \\
\texttt{        But I still don't have the words for 3, 4, 5, 6.} \\
\texttt{        And I don't have the words for road, lake, ship, hill, king, farmer, old, dark, cold.} \\
\texttt{        This is impossible unless I can deduce them.} \\
\texttt{        Is there a pattern in the words?} \\
\texttt{        tree: hobofi} \\
\texttt{        stone: noba} \\
\texttt{        bird: hede} \\
\texttt{        dog: gubo} \\
\texttt{        river: nadu} \\
\texttt{        boat: duha} \\
\texttt{        hunter: feri} \\
\texttt{        spear: nife} \\
\texttt{        big: nafegu} \\
\texttt{        small: bina} \\
\texttt{        tall: hunege} \\
\texttt{        blue: nifafu} \\
\texttt{        red: gabega} \\
\texttt{        find: reho} \\
\texttt{        chase: rinera} \\
\texttt{        hold: fohu} \\
\texttt{        carry: hebo} \\
\texttt{        see: daba} \\
\\ 
\texttt{        Maybe the missing words are formed by analogy?} \\
\texttt{        "road" {-}> ?} \\
\texttt{        "lake" {-}> ?} \\
\texttt{        "ship" {-}> ?} \\
\texttt{        "hill" {-}> ?} \\
\texttt{        "king" {-}> ?} \\
\texttt{        "farmer" {-}> ?} \\
\texttt{        "old" {-}> ?} \\
\texttt{        "dark" {-}> ?} \\
\texttt{        "cold" {-}> ?} \\
\texttt{        "three" {-}> ?} \\
\texttt{        "four" {-}> ?} \\
\texttt{        "five" {-}> ?} \\
\texttt{        "six" {-}> ?} \\
\\ 
\texttt{        Wait, I might have missed a probe translation that contains these words.} \\
\texttt{        Let's check the probe translations again.} \\
\texttt{        1. hobofi reho noba} \\
\texttt{        2. hede rinera gubo} \\
\texttt{        3. hobofi nafegu fohu noba bina} \\
\texttt{        4. nadu nifafu daba duha} \\
\texttt{        5. feri hunege hebo nife gabega} \\
\texttt{        6. hobofiru nafegu ni reho noba bina} \\
\\ 
\texttt{        None of these contain "road", "lake", "ship", "hill", "king", "farmer", "old", "dark", "cold", "three", "four", "five", "six".} \\
\texttt{        This is a contradiction. The exam sentences use words I haven't seen.} \\
\texttt{        Unless... the exam sentences are not using new words?} \\
\texttt{        "the four tall roads" {-}> "road" is new. "four" is new.} \\
\texttt{        "the two birds" {-}> "bird" is new (hede). "two" is new (ni).} \\
\texttt{        "two tall birds" {-}> "tall" (hunege), "birds" (hede).} \\
\texttt{        "four dark hunters" {-}> "dark" is new. "hunters" (feri).} \\
\texttt{        "two red hunters" {-}> "red" (gabega).} \\
\texttt{        "four dark spears" {-}> "spear" (nife).} \\
\texttt{        "two rivers" {-}> "river" (nadu).} \\
\texttt{        "two tall ships" {-}> "ship" is new.} \\
\texttt{        "five blue rivers" {-}> "five" is new.} \\
\texttt{        "three small lakes" {-}> "lake" is new. "three" is new.} \\
\texttt{        "six hills" {-}> "hill" is new. "six" is new.} \\
\texttt{        "cold dog" {-}> "cold" is new.} \\
\\ 
\texttt{        There must be a way to derive these.} \\
\texttt{        Maybe the number words are:} \\
\texttt{        two = ni} \\
\texttt{        three = ?} \\
\texttt{        four = ?} \\
\texttt{        five = ?} \\
\texttt{        six = ?} \\
\texttt{        Maybe they are related to the nouns?} \\
\texttt{        Or maybe I missed a subtle clue.} \\
            }
        }
    }
    & & \\ \\

    \theutterance \stepcounter{utterance}  
    & & & \multicolumn{2}{p{0.3\linewidth}}{
        \cellcolor[rgb]{0.9,0.9,0.9}{
            \makecell[{{p{\linewidth}}}]{
                \texttt{\tiny{[GM$|$GM]}}
                \texttt{expected exactly 10 answer lines, but found 15} \\
            }
        }
    }
    & & \\ \\

    \theutterance \stepcounter{utterance}  
    & & & \multicolumn{4}{p{0.6\linewidth}}{
        \cellcolor[rgb]{0.9,0.9,0.9}{
            \makecell[{{p{\linewidth}}}]{
                \texttt{\tiny{[P1$\langle$GM]}}
                \texttt{Your last response could not be processed: expected exactly 10 answer lines, but found 15.} \\
\texttt{Answer again with exactly 10 numbered lines:} \\
\texttt{ANSWERS:} \\
\texttt{1. <translation>} \\
\texttt{... through 10.} \\
            }
        }
    }
     \\ \\

    \theutterance \stepcounter{utterance}  
    & \multicolumn{4}{p{0.6\linewidth}}{
        \cellcolor[rgb]{0.561,0.737,0.561}{
            \makecell[{{p{\linewidth}}}]{
                \texttt{\tiny{[P1$\rangle$GM]}}
                \texttt{ANSWERS:} \\
\texttt{1. nifafu hunege nifafu reho nifafu hunege ni} \\
\texttt{2. hede ni daba nadu bina} \\
\texttt{3. hede hunege rinera feri hunege ni} \\
\texttt{4. nadu bina daba feri gabega ni} \\
\texttt{5. nadu hunege daba nife gabega ni} \\
\texttt{6. nadu hunege rinera nadu ni} \\
\texttt{7. duha hunege hebo duha nifafu ni} \\
\texttt{8. nadu nifafu daba nadu bina ni} \\
\texttt{9. nife fohu hunege ni} \\
\texttt{10. nadu ni rinera nadu bina} \\
            }
        }
    }
    & & \\ \\

    \theutterance \stepcounter{utterance}  
    & & & \multicolumn{2}{p{0.3\linewidth}}{
        \cellcolor[rgb]{0.9,0.9,0.9}{
            \makecell[{{p{\linewidth}}}]{
                \texttt{\tiny{[GM$|$GM]}}
                {'kind': 'exam', 'correct': 0, 'size': 10, 'results': [0, 0, 0, 0, 0, 0, 0, 0, 0, 0], 'token_similarity': 0.36721334221334223, 'answers': ['nifafu hunege nifafu reho nifafu hunege ni', 'hede ni daba nadu bina', 'hede hunege rinera feri hunege ni', 'nadu bina daba feri gabega ni', 'nadu hunege daba nife gabega ni', 'nadu hunege rinera nadu ni', 'duha hunege hebo duha nifafu ni', 'nadu nifafu daba nadu bina ni', 'nife fohu hunege ni', 'nadu ni rinera nadu bina'], 'targets': ['hageru hunege ne fu reho gibahiru nihe ne fu', 'hederu ni fu fene niri nunu', 'hederu hunege ni rinera feriru hebi ne fu', 'niriru bina ne fu habuga feriru gabega ni', 'niriru hunege ni fene niferu hebi ne', 'hobofi nihe fu rinera naduru ni', 'derahiru hunege ni hebo hederu nifafu ni', 'naduru nifafu dafa daba niriru bina ruga', 'nife fu fohu buriberu hero', 'niriru ni fu rinera gubo nihe fu']}
            }
        }
    }
    & & \\ \\

\end{supertabular}
}

\end{document}
